\section{Archived verification study and limits}
\label{app:empirical}
This appendix incorporates the completed verifier-study evidence into the
current manuscript. It does not restart that study, pool its cohorts with S1,
or turn its numerical positives into source-complete proofs. The seven tables
below are included unchanged from the committed, reproducible summary view
used by the earlier technical-review draft.

\paragraph{Models, cohorts, and costs.}
In configuration names, $L$ is routed depth, $E$ the number of experts per
layer, and $R$ denotes a residual model. Routers have a width-128 ReLU hidden
layer; expert hidden widths are 256 and 128. Sequential networks flatten the
image and use width-128 intermediate routed outputs. Residual models instead
use a linear image-to-128 stem, blocks $h+\operatorname{MoE}(h)$, and a linear
ten-logit head. Both omit batch normalization. Seeds, split, epochs, optimizer,
and checkpoint selection are frozen. Sequential training averages the balance
penalty across layers, while residual training sums it; this is not a causal
ablation of residual connections alone.

The cascade study selects twenty common clean-correct inputs across all
eighteen models from 435 nonexcluded candidates. Its primary comparison uses
the fifteen multilayer models, hence 300 paired model--input objects on only
twenty input blocks. Its conditional positive rate is not S1's whole-prefix
certified accuracy. S1 uses the separate first-200/1000 all-input denominators,
including clean errors and disclosed historical exposure. Neither model nor
input selection may be transferred between these cohorts after outcomes.

Within the historical paired comparison both arms use the same disclosed R7
compatibility snapshot, CPU/double precision, one numerical thread, 8GiB
address-space limit, and 300-second full-request budget. GenBaB uses optimized
alpha bounds and nonlinear branching, with separately charged prefix bounds.
The results describe this configured engine, not every possible configuration.
All 1080 registered B4 calls were audited, with 56.319 charged execution hours;
the 240 unadmitted M-B/M-C ablation slots remain unused. Errors and timeout
costs stay in the denominator. Shared-host variation limits fine timing claims.

\paragraph{What the stopped tests establish.}
The earlier multilayer study has 600 calls: 540 across three configurations,
three seeds, twenty inputs and three arms, plus sixty single-layer pruning
anchors. Pruned GenBaB obtains 108/180 multilayer positives. ACT v1's zero
consists of 178 timeouts, one error and one unsupported result, not 180 solver
timeouts. A supplemental attribution gate registered after a 189-call prefix
finds many-history dominance in 2/52 L3E8 and 0/58 L2E16 objects after removing
replayed unsafe cases. Its timing is disclosed; this is not a fully
preregistered rejection of every possible routing-complexity bottleneck.

M-A replaces gates by independent enclosing intervals, checks last-gate
endpoints, and falls back to coupled solving with the remaining request time
after a phase cap of 120 seconds. Its nine gains all use the relaxed path;
48/52 endpoints close in initial propagation. Three complete-request losses
remain, so net gain is two percentage points with input-block interval [0,4],
below the registered five-point threshold. Gains concentrate in two-layer
models. The paid repeat of the coupled arm changes four of 120 terminal states,
including one positive-to-timeout change: that is not a three-percentage-point
estimate of positive-result noise. Neither a shared positive nor a retained
candidate refunds time spent on an unsuccessful cascade phase.

M-B uses nonnegative guard multipliers and weak duality. The original
independent-trajectory comparison could not isolate incremental guard benefit.
The revised test retains an evaluated twenty-step alpha prefix, then compares
forty guard/alpha updates against paid extra-alpha updates on the same eleven
histories. It adds one complete local closure, below the required five.
Original M-C partitions the first-layer margin of twenty-seven local histories
and adds no closures. A proposed last-layer dual-conditioned variant was gated
on revised M-B and was not executed. These outcomes do not prove that every
guard, margin partition or expert-network representation is ineffective.

The larger-radius and single-layer results below do not overturn those stops.
No gain appears among the thirteen larger-radius objects with multiple observed
histories; a single observed history elsewhere does not prove route stability.
The single-layer cohort has the same 57 positives in both arms, with extra
cascade cost. None of these diagnostic observations isolates ReLU difficulty
as the unique cause of every unresolved request.

% Relevant committed summary primary sha256 30c1121e079619babae3fb0bd0d6f09652d1e1be3ca0efb126b35342b94e875d
% Relevant committed summary multilayer sha256 5e78f74054d0dd883fbfe692c8a5146f155786fb917f0c491c54be8a49c6d127
% Relevant committed summary guard sha256 8bebc40ef4f01d9630a61de608044f2d730faf7b2574acf165b5262afa56f8fa
\begin{table}[ht]
\centering\small
\caption{Complete-request cascade: 300 calls per arm. P is numerical POSITIVE, not a strict certificate. All terminals and costs are retained.}
\label{tab:primary}
\begin{tabular}{lrrrrrr}
\toprule
Method & P & UNKNOWN & TIMEOUT & ERROR & Mean s & Median s \\
\midrule
Coupled & 116 & 8 & 176 & 0 & 189.287 & 288.350 \\
Cascade & 122 & 5 & 171 & 2 & 192.983 & 288.337 \\
\bottomrule
\end{tabular}
\end{table}

\begin{table}[ht]
\centering\small
\caption{Family breakdown of the same primary cohort: 60 model--input objects per row. Gain/loss counts are paired, not differences between unmatched totals.}
\label{tab:families}
\begin{tabular}{lrrrr}
\toprule
Family & Coupled P & Cascade P & Gained & Lost \\
\midrule
L2E8 & 24 & 28 & 5 & 1 \\
L3E8 & 27 & 26 & 0 & 1 \\
L2E16 & 53 & 53 & 0 & 0 \\
R2E8 & 7 & 10 & 4 & 1 \\
R3E8 & 5 & 5 & 0 & 0 \\
\bottomrule
\end{tabular}
\end{table}

\begin{table}[ht]
\centering\small
\caption{Separate multilayer bottleneck study: 60 objects per configuration/arm. The single-layer cost anchor is a separate pruning-only 60 calls.}
\label{tab:multilayer}
\begin{tabular}{lrrr}
\toprule
Configuration & ACT v1 P & Naive P & Pruned P \\
\midrule
L2E8 & 0 & 0 & 28 \\
L3E8 & 0 & 0 & 33 \\
L2E16 & 0 & 0 & 47 \\
\bottomrule
\end{tabular}
\end{table}

\begin{table}[ht]
\centering\small
\caption{Revised guard diagnostic on the same eleven histories. Each continuation retains its complete twenty-evaluation baseline. A local history is not a full MoE request.}
\label{tab:guard}
\begin{tabular}{lrr}
\toprule
Continuation & Closed histories & UNKNOWN \\
\midrule
Extra alpha & 0 & 11 \\
Joint alpha/guard & 1 & 10 \\
\bottomrule
\end{tabular}
\end{table}

\clearpage

\section{S1 seed-level counts and evidence reconstruction}
\label{app:s1seeds}
Table~\ref{tab:s1seeds} preserves all twelve fixed models, including the unmatched
L2 extension. Complete and incomplete columns have different denominators and
must not be added. Every $8/255$ positive count is zero. Three seeds do not
constitute a sample of all architectures, and the input-block intervals in the
main text do not quantify random-training uncertainty.

% Seed-count table inserted from the committed S1 final derived archive.
% S1 archive sha256 8dde94094d1fe97068ffa55dd32dd6204a8a9a911e089f292aa77c5f962adec5
\begin{table}[ht]
\centering\small
\caption{S1 numerical positives by fixed seed. C denotes complete verification (out of 200); I denotes incomplete CROWN-IBP (out of 1000 per model). No seed is omitted.}
\label{tab:s1seeds}
\begin{tabular}{llrrrr}
\toprule
Group & Seed & C ($2/255$) & I ($2/255$) & C ($8/255$) & I ($8/255$) \\
\midrule
L1E8 & 0 & 46 & 16 & 0 & 0 \\
L1E8 & 1 & 34 & 3 & 0 & 0 \\
L1E8 & 2 & 34 & 4 & 0 & 0 \\
L2E8 & 0 & 4 & 0 & 0 & 0 \\
L2E8 & 1 & 7 & 0 & 0 & 0 \\
L2E8 & 2 & 6 & 0 & 0 & 0 \\
Dense total & 0 & 26 & 0 & 0 & 0 \\
Dense total & 1 & 47 & 1 & 0 & 0 \\
Dense total & 2 & 42 & 1 & 0 & 0 \\
Dense active & 0 & 40 & 1 & 0 & 0 \\
Dense active & 1 & 35 & 1 & 0 & 0 \\
Dense active & 2 & 41 & 0 & 0 & 0 \\
\bottomrule
\end{tabular}
\end{table}

The committed S1 archive is
\begin{quote}\small
\path{docs/architecture_s1_result_20261031.json}.
\end{quote}
It contains 48 model/mode/radius rows, all eight paired comparisons, gained/lost input-seed
identities, and costs. Its audit hash binds the full original terminal records.
The eight old public-PGD errors are retained, accounting for eighty unaccepted
domains; they are neither rerun nor imputed as successful attacks. All 4800
complete and 24000 incomplete verification requests are settled without missing
costs. Shared-clean-correct strata explain all-request timing but do not alter
the registered primary denominator. S2 remains unexecuted after the stop.

The existing table script is
\begin{quote}\small
\path{scripts/render_current_review_tables.py}.
\end{quote}
Running it with \texttt{python -I -S -B} and \texttt{--check} reconstructs the
seven historical tables from their committed source summaries.
It does not load a model, execute an attack, or independently verify a numerical
bound. The current manuscript build separately checks the S1 transcription
against its final archive. Table reconstruction is not clean-environment
empirical reproduction; a human reviewer must distinguish those tasks.

\section{Source contracts and independent proof boundary}
\label{app:sourceboundary}
The historical output-layer ACT study reports 179 policy acceptances against
156 for matched monolithic and 141 for the older strong path, on 100 inputs
and three models. A later audit identifies input-containment gaps affecting
all 23 primary gains. Those counts remain numerical-policy outcomes; they are
not source-complete certificates. The requested and reconstructed input sets
are not nested. This does not establish that the models are unsafe, or predict
the results of a corrected execution. No sealed case is reopened here.

For a supplied LP $Ax\leq b$, $l\leq x\leq u$, objective $c^\top x+c_0$, and
any $y\leq0$, set $r=c-A^\top y$. Exact arithmetic can check the bound
\[
 c_0+b^\top y+\sum_i\min\{r_i l_i,r_i u_i\}.
\]
Indeed $y^\top Ax\geq y^\top b$, and the residual term is minimized over the
variable box. Multiplier signs, original coefficients, residual compensation,
property identity and obligation coverage must all be checked. This proves a
statement about the supplied LP. It does not establish input inclusion,
network-to-HZ conversion, guard lowering, omitted-history infeasibility or
native floating-point execution. A positive bound from an old matrix cannot
be attached to a new checked source construction.

Human technical review, a clean empirical reproduction, authorship and any
distribution and submission decisions remain external tasks. The manuscript and
record audits do not fill them or certify that the contribution meets a
particular venue's novelty standard. All unsuccessful gates, original errors,
and source limitations remain visible alongside local positive results.
